\documentclass[12pt]{article}
\usepackage[utf8]{inputenc}
\usepackage[T1]{fontenc}
\usepackage{hyperref}
\title{Northern Amami-Oshima Language Resources for K-12 Curriculum Development}
\author{Compiled by Claude Sonnet 5.5 (claude-sonnet-5-5) for Andr\'as Kornai}
\date{}
\begin{document}
\maketitle

\section{Introduction}
Northern Amami-Oshima (local name shimayumuta) is a Ryukyuan language of the north of Amami Oshima, Kagoshima Prefecture, Japan. It is endangered. Wikipedia gives about 12,000 native speakers for 2004. The results contained no school curriculum for it. This dossier is short because the search found little beyond reference pages and one research record.

\section{Primary Educational Resources}
No textbooks or school programmes were found. A Kagoshima University paper surveys shimauta (island song) classes on Amami Oshima, including community-centre courses: \url{https://ir.kagoshima-u.ac.jp/records/4917}. Another Kagoshima University record notes that shimauta competitions and shimauta schools have multiplied: \url{https://ir.kagoshima-u.ac.jp/records/10480/export/json}. Neither abstract was read in full.

\section{Linguistic and Reference Materials}
\begin{itemize}
\item Wikipedia, Amami \=Oshima language: \url{https://en.wikipedia.org/wiki/Amami_%C5%8Cshima_language}
\item Omniglot, Amami (\emph{shimayumuta}): \url{https://omniglot.com/writing/amami.htm}
\item Japanese Wiktionary page for the language: \url{https://ja.wiktionary.org/wiki/%E5%8C%97%E5%A5%84%E7%BE%8E%E5%A4%A7%E5%B3%B6%E8%AA%9E}. Only navigation text was retrieved.
\item LINGUIST List issue 24.388: \url{https://dev.linguistlist.org/issues/24/388/}. Its relevance was not checked.
\end{itemize}

\section{Religious and Community Resources}
Shimauta is a song genre of the Amami Islands that is tied to community identity. Overview: \url{https://www.wikipedia.org/wiki/Shimauta}. The Joshua Project page (\url{https://joshuaproject.net/people_groups/14009}) was returned, but its content may concern a different group and should be treated with caution.

\section{Digital and Audio Resources}
No dedicated audio corpus or software was found. Shimauta recordings would be a natural source, but none were found in the results.

\section{Curriculum Development Context}
Shimauta classes show community interest in the island language and could be a starting point for school-level work. No existing dictionary or textbook was identified in these results. Further searches for Japanese-language dictionaries and local education-board materials are needed.

\end{document}